% STATUS: written 2026-11-06 (English gate clean; sources read in full as recorded below).
% Sources read in full: Kadison--Ringrose 1971 (Acta Math. 126, 227--243). Johnson 1972 not secured (no open PDF); marked below.
\section{Kadison--Ringrose and bounded cohomology}\label{sec:cohomology:bounded-cohomology}

For operator algebras the continuous structure
matters~\cite[pp.~227--228]{KadisonRingrose1971}. Kadison and Ringrose admit
only bounded maps as cochains, and the bound is what changes the
theory~\cite[pp.~227--228]{KadisonRingrose1971}. These maps give cohomology
groups, written
\[
  H_w^n(\mathcal{R}, \mathcal{M}),
\]
and the letter $w$ stands for weak~\cite[pp.~227--228]{KadisonRingrose1971}.

Their first theorem of this kind is a vanishing statement, and it comes
next~\cite[Theorem~4.4, p.~239]{KadisonRingrose1971}. On the type I von Neumann
algebras the whole bounded theory
vanishes~\cite[Theorem~4.4, p.~239]{KadisonRingrose1971}.

\begin{theorem}[{\cite[introduction and Theorem 4.4, pp.~227--228, 239]{KadisonRingrose1971}}]\label{thm:cohomology:kadison-ringrose}
$H_w^n(\mathcal{R}, \mathcal{R}) = 0$ for every $n \ge 1$ when $\mathcal{R}$ is a
type I von Neumann algebra.
\end{theorem}

In degree one this is an older
statement~\cite[pp.~227--228]{KadisonRingrose1971}. A von Neumann algebra has no
outer derivations~\cite[pp.~227--228]{KadisonRingrose1971}. So every derivation
of it is inner~\cite[pp.~227--228]{KadisonRingrose1971}. Read through the
Hochschild theory, that says the first group vanishes
\cite[pp.~227--228]{KadisonRingrose1971}. So
Theorem~\ref{thm:cohomology:kadison-ringrose} extends the vanishing of
derivations to every degree.

% STATUS: written 2026-11-06 with one source unobtainable. Johnson 1972 is an
% AMS Memoir with no open copy and no shadow-library copy: the searches run on
% 2026-11-06 through papergraph found only the unrelated 1975 "Introduction to
% cohomology in Banach algebras", which is a different work. The section names
% Johnson's theory and points to the memoir in the bibliography without quoting
% a statement from it. Recorded in ledger/manual-downloads.md.
The theory of Johnson for amenable algebras is the bounded analogue of the
separability of Hochschild~\cite[pp.~227--228]{KadisonRingrose1971}.
The memoir that reports it is recorded in the bibliography~\cite{Johnson1972}.
No copy of that memoir could be secured, so no statement of it is quoted here.
A reader who has the memoir will find the bounded cohomology of an amenable
algebra treated there.
